\chapter{Introduction}
\label{chap:introduction}

This blueprint describes the mathematics formalized in the Lean library gc-zk: zero-knowledge
proofs built from privacy-free garbled circuits. Every definition and theorem below names its Lean
counterpart, and the web version of this document has a dependency graph that shows how the
results build on each other. All the statements in this document
are proved in Lean, with no \texttt{sorry}, and the main theorems depend only on the standard
axioms of Lean.

\section*{Zero knowledge from garbled circuits}

The idea of a garbled circuit is due to Yao, who presented it orally in the context of secure
two-party computation; the construction does not appear in his written paper \cite{Yao86}, and its
history is told by Bellare, Hoang and Rogaway \cite[Sect.~1]{BHR12b}. In 2013 Jawurek,
Kerschbaum and Orlandi \cite{JKO13} observed that they also give an efficient zero-knowledge proof
for any statement of the form ``I know $w$ with $C(w) = 1$'' for a Boolean circuit $C$. The roles
are reversed: the verifier garbles $C$, the prover evaluates it on its witness, and the garbling
is opened at the end so that the prover can check that it was honest. Since the prover knows its
own input, the garbling does not need to hide anything from the evaluator, only to prevent it
from forging output labels. Frederiksen, Nielsen and Orlandi \cite{FNO15} called such schemes
\emph{privacy-free} and showed that they are considerably cheaper. With Free XOR \cite{KS08} and
the half-gates of Zahur, Rosulek and Evans, XOR gates cost nothing and AND gates cost one
ciphertext \cite[Sect.~6]{ZRE15}.

The garbling scheme formalized here is the privacy-free scheme of Zahur, Rosulek and Evans
\cite[Sect.~6]{ZRE15}: Free XOR, and for the non-linear gate their one-ciphertext evaluator
half-gate \cite[Sect.~3.1]{ZRE15}, written as a multiplexer (Remark~\ref{rem:mux-half-gate}).
Circuits are built from inputs, constants, XOR gates and multiplexers.

\section*{What the protocol is for}

The protocol is an interactive zero-knowledge argument of knowledge (soundness holds against
provers with bounded resources) for statements of the form ``I
know $w$ with $C(w) = 1$'', where $C$ is any Boolean circuit known to both parties. Its cost is
proportional to the size of $C$ \cite[Sect.~1]{JKO13}, and it uses only symmetric primitives (a hash) besides the
oblivious transfers, one per bit of the witness. It is therefore suited to statements without
algebraic structure, for which the classical zero-knowledge protocols for discrete logarithms or
other group relations do not apply.

\emph{Typical uses.} The examples of Jawurek, Kerschbaum and Orlandi \cite[Sect.~1 and~1.1]{JKO13}
are of this kind:
\begin{itemize}
  \item knowledge of a preimage, ``I know $x$ with $y = \mathrm{SHA\text{-}256}(x)$'' for a public
  $y$;
  \item knowledge of a key, ``I know $k$ with $y_1 = \mathrm{AES}_k(y_2)$'';
  \item statements about data encrypted with a block cipher, for instance that part of a
  ciphertext decrypts to a given value, without releasing the key;
  \item knowledge of a signature, or of a Fiat--Shamir proof, without revealing it, which requires
  proving statements about the hash function that instantiates the random oracle.
\end{itemize}
They report experiments for AES, MD5 and SHA-256 \cite[Sect.~5.2]{JKO13}. With the privacy-free
half-gates garbling formalized here, the garbled circuit for ``I know $k$ such that
$\mathrm{AES}(k, m) = c$'' takes about 108 kilobytes with 128-bit labels \cite[Sect.~6]{ZRE15}.

\emph{What it does not give.}
\begin{itemize}
  \item \emph{It is interactive, and the proof convinces only the verifier.} The verifier's
  messages are not just random challenges, so the protocol is not public-coin and the Fiat--Shamir
  transformation does not apply \cite[Sect.~3, footnote~6]{JKO13}. The protocol as formalized here
  has no non-interactive or publicly verifiable version: the verifier itself builds the garbled circuit and later opens it,
  so a transcript is not evidence for anyone else. It does not fit settings where a proof must be
  posted once and checked by many.
  \item \emph{One garbling per proof.} At the end the verifier opens its seed, which reveals the
  global offset $\Delta$, so a garbled circuit cannot be reused for another proof.
\end{itemize}

\emph{Ways to make it insecure.} Each of the following breaks a hypothesis of the theorems below,
or makes their bounds useless.
\begin{itemize}
  \item \emph{Skipping the prover's check.} A malicious verifier can garble a circuit other than
  $C$, or mount a selective-failure attack on the input phase \cite[Sect.~1]{JKO13}. Zero knowledge
  rests on the prover committing to its output label, verifying the opened garbling against $C$
  and the labels it received, and only then revealing the label (Theorem~\ref{thm:zk}). A prover
  that reveals the label before the check, or that does not check, leaks information about its
  witness; Proposition~\ref{prop:no-check-leaks} gives an explicit example.
  \item \emph{Sending the garbled circuit before the oblivious transfers, or using an OT that lets
  the prover choose its input late.} If the circuit is sent before the oblivious transfers, the
  garbling must be secure against adaptive attacks \cite[Sect.~3, footnote~7]{JKO13}. The same
  holds whenever the prover can choose its witness bits after seeing the garbled circuit.
  This is the situation of the Chou--Orlandi OT, where a receiver can delay its decryption; it is
  treated in Chapters~\ref{chap:adaptive} and~\ref{chap:co-soundness}.
  \item \emph{Weak oblivious transfer keys.} With the PVW OT, a receiver key $(0, 0)$ makes every
  ciphertext carry its message in the clear; the sender in this library rejects every key whose
  first component is $0$ (Lemma~\ref{lem:pvw-ext}). With the Chou--Orlandi OT, a secret $a_j = 0$ makes the two masks of
  input $j$ equal; this costs $k/q$ in the soundness bound (Theorem~\ref{thm:co-sound}).
  \item \emph{An unsuitable hash.} Correlation robustness of the hash alone is not enough to prove
  Free XOR secure: Choi, Katz, Kumaresan and Zhou rule out a fully black-box reduction
  \cite[Sect.~3]{CKKZ12}. The theorems here assume a random oracle; for Free XOR in Yao's protocol
  a circular form of correlation robustness suffices \cite[Sect.~4]{CKKZ12}.
  \item \emph{Small parameters.} The bounds below are of order $q/|V|$ for $q$ hash queries; the
  labels must be long enough (for instance 128 bits) for this to be negligible, and likewise the
  commitment values and, with a real oblivious transfer, the order of the group.
\end{itemize}

\section*{No claim of originality}

No originality is claimed for any definition, statement or proof idea in this library or in this
document. The protocols, the security notions and the arguments come from the cited literature, or
are standard in it; what is ours is only their formalization in Lean. Where we could not yet match a
statement with a published one, we say so. Should any of them turn out to be new, this will be
stated separately.

\section*{Contents}

Chapter~\ref{chap:variants} explains how the order of the messages (when the witness is fixed, when
the garbling is opened) and the choice of commitment and oblivious transfer determine the security
that can be proved, discusses security against adaptive corruptions, and points to the theorems
that cover each case. Chapter~\ref{chap:garbling} defines circuits, the garbling and the evaluation, and proves
correctness. Chapter~\ref{chap:jko} formalizes the protocol of \cite{JKO13} in the hybrid model
with ideal oblivious transfer and commitment, and proves completeness, perfect zero knowledge
against a malicious verifier, and soundness by reduction to authenticity.
Chapter~\ref{chap:authenticity} shows that $\Delta$ reaches the evaluator only through the hashes
of inactive labels, and Chapter~\ref{chap:random-oracle} bounds the authenticity advantage of a
$q$-query adversary in the random oracle model by $(q+1)/(|V|-1)$, for the real garbling with
uniform zero-labels. Chapter~\ref{chap:commitment} replaces the ideal commitment by a hash
commitment in the random oracle model, with statistical zero knowledge and soundness by
extraction. Chapter~\ref{chap:ot} replaces the ideal OT by the Chou--Orlandi OT and proves zero
knowledge with both real components. Chapter~\ref{chap:pvw} uses instead the dual-mode OT of
Peikert, Vaikuntanathan and Waters \cite{PVW08}, which realizes the OT functionality of
\cite{JKO13}: soundness holds in messy mode, and zero knowledge is reduced to DDH by explicit
distinguishers. With the Chou--Orlandi OT the prover can choose its input after seeing the garbled
circuit, so soundness needs \emph{adaptive} authenticity: Chapter~\ref{chap:adaptive} proves that
an adversary that interleaves hash queries and label requests wins with probability at most
$(q+1)/(|V|-1) + q/|V|$, and Chapter~\ref{chap:co-soundness} uses it to prove soundness with the
Chou--Orlandi OT, by reduction to the square Diffie--Hellman problem.
Chapter~\ref{chap:cut-and-choose} contains the cut-and-choose bounds against a malicious garbler,
derived from the garbling, and Chapter~\ref{chap:limitations} lists what is not formalized.

\section*{The model}

All the security statements are information-theoretic statements in idealized models (the random
oracle model \cite{BR93}, and the hybrid models of \cite{JKO13}), or reductions to explicit
adversaries against DDH and square Diffie--Hellman. They are proved by counting: a probability is
the proportion of the randomness on which an event holds, and the proofs are sequences of games
\cite{Sho04} in which ``identical until bad'' \cite{BR06} becomes an injection between sets of
randomness. Adversaries are deterministic; a randomized adversary is a family of them indexed by
its coins, and every bound holds for each member of the family, hence on average. The numbers of
queries are fixed; an adversary that makes fewer queries, or interleaves its queries to several
oracles differently, is brought to this form by padding with dummy queries, an argument that is
stated but not formalized. Lean has no notion of efficient computation, so polynomial-time
constraints, and the efficiency of the reductions, are outside the formalization.

\emph{Notation.} Labels live in a vector space $V$ over $\F$, and $\Delta \in V$ is the global
offset. The letter $q$ has two uses: in Chapters~\ref{chap:random-oracle} and~\ref{chap:adaptive}
it is the number of queries (or actions) of the adversary; in Chapters~\ref{chap:ot},
\ref{chap:pvw} and~\ref{chap:co-soundness} it is the order of the group $\mathbb{G}$ with
generator $g$, and numbers of queries are written $q_0, \dots, q_3$, $Q$ or $n_0, n_1, n_2$.
Chapter~\ref{chap:variants}, which summarizes all of them, uses $q$ in both senses, each time as
in the chapter it refers to. From
Chapter~\ref{chap:commitment} on, $G$ is the commitment oracle, distinct from the group
$\mathbb{G}$. A
circuit input is $x$ (or the witness $w$), and $b$ is a bit, except in the oblivious transfers,
where $b$ is the receiver's secret exponent. In Chapter~\ref{chap:pvw}, $b$ is also the branch of
the cryptosystem, and $x_0$, $x_1$, $x$, $y$ are exponents of the common string.

\section*{Related work}

\emph{Formalized cryptography.} Game-based cryptographic proofs have been mechanized in several
frameworks, among them EasyCrypt (see \cite{ABB+17}), CryptHOL in Isabelle/HOL \cite{BLS20},
SSProve in Coq \cite{HRV+23}, and, in Lean~4, VCVio \cite{TDW+26} and CatCrypt \cite{Spi26}. Closest to the present library are the EasyCrypt verification of
Yao's secure function evaluation protocol, with a formalization of garbled circuits in the framework
of Bellare, Hoang and Rogaway \cite{ABB+17, BHR12b}, and the formalization of simulation-based
security of multi-party computation in Isabelle/HOL by Butler, Aspinall and Gascón \cite{BAG17}.
In Lean~4, Dziembowski, Fabia\'nski, Micciancio and Stefa\'nski formalize the computational
soundness of a symbolic framework and, as an extended example, the simulation-based privacy of a
garbling scheme taken from Li and Micciancio, proved symbolically \cite[Sect.~3]{DFMS26}.
Among its 172 protocols, CatCrypt lists Yao's garbled circuits with half-gates, row reduction and
Free XOR, and a DDH-based ``Simplest OT'' \cite[Sect.~4]{Spi26}; its protocol sources are in a
private repository \cite[Sect.~7]{Spi26} (only the core is public), so we have not compared its definitions or bounds with ours. As of
November 2026, the README of the public core also lists machinery for universal composability: UC
emulation, the dummy-adversary theorem and a composition law, and says that the
universal-composability theory itself is in separate distributions \cite{CCcore}.
We are not aware of a previous formalization of zero knowledge from garbled circuits or of
privacy-free garbling. The present library does not use a general framework: its probabilities are counts over
finite types, and its games are written directly as functions of the randomness.

\emph{Adaptive garbling.} Adaptive security of garbling, where the input is chosen after the
garbled circuit, was introduced by Bellare, Hoang and Rogaway \cite{BHR12}, who define adaptive
privacy \cite[Sect.~3, Fig.~2]{BHR12} and adaptive obliviousness and authenticity
\cite[Sect.~4, Fig.~6]{BHR12}, and give transforms that
make any statically secure scheme adaptively secure. Hemenway, Jafargholi, Ostrovsky, Scafuro and
Wichs prove simulation-based adaptive security, a privacy notion, for a modification of Yao's scheme
under one-way functions \cite[Sect.~3 and~5]{HJO+16}; Jafargholi and Wichs prove the same notion for
Yao's scheme itself, unmodified, with a security loss governed by a pebbling game
\cite[Sect.~3 and Theorem~1]{JW16}. Adaptive security of unmodified half-gates was studied in idealized models by Barnum,
Heath, Kolesnikov and Ostrovsky \cite{BHKO26} and by Guo, Yang, Wang, Yu and Liu \cite{GYWYL}.
Chapter~\ref{chap:adaptive} needs adaptive authenticity of unmodified privacy-free garbling, for
which we have not found a published statement.
Ganesh, Kondi, Patra and Sarkar \cite[Sect.~3]{GKPS18} show that the protocol of \cite{JKO13} is secure
against \emph{adaptive corruptions} when the oblivious transfer satisfies a mild adaptive security
guarantee; this is a different notion from the adaptive choice of the input considered here.

\emph{The Chou--Orlandi OT.} The issue that motivates Chapters~\ref{chap:adaptive}
and~\ref{chap:co-soundness}, that a receiver can delay the decryption, was pointed out by Byali,
Patra, Ravi and Sarkar \cite[Sect.~3]{BPRS17} and discussed by Doerner, Kondi, Lee and shelat
\cite[App.~A]{DKLs18}, who add a verification step to the protocol; Chou and Orlandi refer to both
\cite[Sect.~1.1]{CO15} and withdraw the claim of UC security of their protocol.

\section*{How this library was written}

The author chose the protocol, the model and the statements, and reviewed the results. The Lean
proofs were written by the AI assistant Claude (Anthropic) working under his direction, and so was
a first draft of the pen-and-paper arguments of Chapters~\ref{chap:adaptive}
and~\ref{chap:co-soundness} (the linear parts of labels, the bad steps, and the reduction with
guessed key queries), which the author then reviewed. Correctness does not depend on trusting
either of them: Lean checks every proof. The formalization uses Lean~4 \cite{MU21} and
Mathlib \cite{Mathlib20}, and this document is built with leanblueprint \cite{Mas}.
